\clearpage
\appendix
\section{Código completo}
\label{ap:codigo}

Este apêndice traz o código numérico completo usado na análise. Os scripts simbólicos
da Seção~\ref{sec:intro} já aparecem por inteiro nos
Códigos~\ref{lst:solucao_geral} a~\ref{lst:frf} e não são repetidos.

Os arquivos ficam na pasta \texttt{codigo/} e leem os dados de \texttt{dados/}. Para
reproduzir as figuras e tabelas, basta executar, dentro de \texttt{codigo/},
\begin{center}
    \texttt{python gera\_resultados.py}
\end{center}
com Python 3.11 ou mais recente e os pacotes listados em \texttt{requirements.txt}
(NumPy, SciPy, Matplotlib, SciencePlots e SymPy). O estilo SciencePlots usa o
\LaTeX{} do sistema para escrever os textos das figuras. Os comentários
\texttt{\# <{}<nome>{}>} marcam os trechos reproduzidos ao longo do texto.

\lstinputlisting[caption={\texttt{parametros.py}: parâmetros da viga e conversão
            $\lambda_r \leftrightarrow f_r$.}, label={lst:ap_parametros}]{codigo/parametros.py}

\lstinputlisting[caption={\texttt{dados.py}: leitura dos arquivos e FRF do martelo a
            partir da série temporal.}, label={lst:ap_dados}]{codigo/dados.py}

\lstinputlisting[caption={\texttt{identificacao.py}: picos e método da meia
            potência.}, label={lst:ap_identificacao}]{codigo/identificacao.py}

\lstinputlisting[caption={\texttt{modelo.py}: equações características, modos e
            FRF analítica.}, label={lst:ap_modelo}]{codigo/modelo.py}

\lstinputlisting[caption={\texttt{ajuste\_molas.py}: ajuste das rigidezes da base,
            validação e sensibilidade.}, label={lst:ap_ajuste}]{codigo/ajuste_molas.py}

\lstinputlisting[caption={\texttt{gera\_resultados.py}: processamento completo,
            figuras e tabelas.}, label={lst:ap_gera}]{codigo/gera_resultados.py}
